\section{Physics-Guided Simulation Pipeline}\label{sec:generation}

Modern generators render physically plausible motion, yet prompt-driven synthesis supplies no known
physical target: because physics and appearance are produced jointly, the intended trajectory,
contact, or counterfactual is neither fixed in advance nor recoverable afterward, and every physical
property stays confounded with appearance. An evaluation testbed instead needs a target factored out
of appearance---specified, known by construction, and auditable independently of generator
capability. We obtain one from a trusted physics solver and render it as a geometric control, so the
solver fixes the motion while the prompt supplies only appearance (\cref{fig:genpipe}). Two stages
develop this: \cref{sec:gen-control} turns the event description into a validated target rendered as a
registered control, and \cref{sec:gen-mechanism} conditions the generator on that control and an
appearance prompt. Whether the generated video realizes the target is the open question of
\cref{sec:gen-experiments}.

\begin{figure}[t]
  \centering
  % =====================================================================
% fig_authoring_pipeline.tex -- how a physics description in text becomes the
% control video the diffusion backbone follows: a language model writes the
% scene specification, the specification is simulated and validated on the
% trajectory it actually produces, a failing specification is returned to the
% model with the violated condition attached and re-authored, and only an
% accepted specification is rendered.
%
% \input-able into main.tex, which loads adobe-techreport.sty (hence the
% adobered / adobeink / adobegray colours and the Source font family).
%
% REQUIRED IN THE MAIN PREAMBLE:
%     \usepackage{tikz}
%     \usetikzlibrary{arrows.meta,positioning,calc}
% Both lines are required, and the libraries MUST be loaded in the preamble --
% see the guard below for why a figure file cannot load them for itself.
% Every style is declared inside the tikzpicture's own option list, so all
% definitions are local to this picture: \input-ing this file twice, or
% alongside its companion figures, cannot clash.
%
% CONVENTIONS shared with fig_critic_pipeline.tex and fig_graph_construction.tex:
%   adobegray!12 fill      = a learned component
%   white + adobeink line  = a deterministic component
%   adobered line          = the one check the pipeline turns on (used here for
%                            validation of the simulated trajectory, and for
%                            nothing else)
%   dashed gray            = something handed to the system, not produced by it
% Every stage in this picture is implemented, so every line is solid; the
% solid/dashed built-versus-roadmap distinction used elsewhere in the report
% does not apply here.
%
% WRAPPING: like every other figs/*.tex in this report, this file contains no
% figure environment and no \caption -- the caller supplies both:
%     \begin{figure}[tp]
%       \centering
%       % =====================================================================
% fig_authoring_pipeline.tex -- how a physics description in text becomes the
% control video the diffusion backbone follows: a language model writes the
% scene specification, the specification is simulated and validated on the
% trajectory it actually produces, a failing specification is returned to the
% model with the violated condition attached and re-authored, and only an
% accepted specification is rendered.
%
% \input-able into main.tex, which loads adobe-techreport.sty (hence the
% adobered / adobeink / adobegray colours and the Source font family).
%
% REQUIRED IN THE MAIN PREAMBLE:
%     \usepackage{tikz}
%     \usetikzlibrary{arrows.meta,positioning,calc}
% Both lines are required, and the libraries MUST be loaded in the preamble --
% see the guard below for why a figure file cannot load them for itself.
% Every style is declared inside the tikzpicture's own option list, so all
% definitions are local to this picture: \input-ing this file twice, or
% alongside its companion figures, cannot clash.
%
% CONVENTIONS shared with fig_critic_pipeline.tex and fig_graph_construction.tex:
%   adobegray!12 fill      = a learned component
%   white + adobeink line  = a deterministic component
%   adobered line          = the one check the pipeline turns on (used here for
%                            validation of the simulated trajectory, and for
%                            nothing else)
%   dashed gray            = something handed to the system, not produced by it
% Every stage in this picture is implemented, so every line is solid; the
% solid/dashed built-versus-roadmap distinction used elsewhere in the report
% does not apply here.
%
% WRAPPING: like every other figs/*.tex in this report, this file contains no
% figure environment and no \caption -- the caller supplies both:
%     \begin{figure}[tp]
%       \centering
%       % =====================================================================
% fig_authoring_pipeline.tex -- how a physics description in text becomes the
% control video the diffusion backbone follows: a language model writes the
% scene specification, the specification is simulated and validated on the
% trajectory it actually produces, a failing specification is returned to the
% model with the violated condition attached and re-authored, and only an
% accepted specification is rendered.
%
% \input-able into main.tex, which loads adobe-techreport.sty (hence the
% adobered / adobeink / adobegray colours and the Source font family).
%
% REQUIRED IN THE MAIN PREAMBLE:
%     \usepackage{tikz}
%     \usetikzlibrary{arrows.meta,positioning,calc}
% Both lines are required, and the libraries MUST be loaded in the preamble --
% see the guard below for why a figure file cannot load them for itself.
% Every style is declared inside the tikzpicture's own option list, so all
% definitions are local to this picture: \input-ing this file twice, or
% alongside its companion figures, cannot clash.
%
% CONVENTIONS shared with fig_critic_pipeline.tex and fig_graph_construction.tex:
%   adobegray!12 fill      = a learned component
%   white + adobeink line  = a deterministic component
%   adobered line          = the one check the pipeline turns on (used here for
%                            validation of the simulated trajectory, and for
%                            nothing else)
%   dashed gray            = something handed to the system, not produced by it
% Every stage in this picture is implemented, so every line is solid; the
% solid/dashed built-versus-roadmap distinction used elsewhere in the report
% does not apply here.
%
% WRAPPING: like every other figs/*.tex in this report, this file contains no
% figure environment and no \caption -- the caller supplies both:
%     \begin{figure}[tp]
%       \centering
%       \input{figs/fig_authoring_pipeline}
%       \caption{...}
%     \end{figure}
% The \label is carried HERE (see the block just below \makeatother), so main.tex
% must NOT add a second \label{fig:authoring}.
%
% CONTENT follows the subsection "From a physics description to a control video"
% (\ref{sec:gen-authoring}) and is drawn from the implementation in
% ~/diffphy_exp014/scripts/exp014/agentic_mask.py -- stage order from build_one()
% at lines 589-663, with:
%   text-in, no reference video   docstring l.4-8; author() l.294-306; main() l.683
%   scene specification written   author() l.294-327; scene contract l.124-149
%   required fields re-asked once _validate() l.249-265; author() l.322-326
%   waypoints only where needed   docstring l.26-29; motion model l.98-117
%   simulated by the renderer's   _sim_traj() l.348-352 (expand + simulate);
%     own engine                  silhouette_render.render_frames l.209 calls the
%                                 same simulate(); physics_world.py is pymunk
%   validation on the trajectory  lint_scene() l.355-450 (ballistic must be
%                                 dynamic l.406; off-frame l.410-417; rest /
%                                 support l.419-424; rising must rise l.426-428;
%                                 reach / miss targets l.430-445)
%   repair loop, three rounds     build_one() l.604-612; repair() l.472-488;
%                                 realized-motion readout _traj_summary l.453-469
%   appearance written apart      author_texture() l.272-291; called l.616
%   camera + graphic-body strip   apply_camera() l.491-504; _strip_graphic_bodies
%                                 l.210-215; called l.618-619
%   render + contact sheet        build_one l.623-627; render_gif; strip l.507-521
%   what the generator receives   build_one l.632-662 (fill-ratio cap l.632-638,
%                                 low-control-strength cap l.652-654), consumed by
%                                 generate_worker.py l.176-183 -> video_stylizer
%                                 WanVACEStylizer.stylize l.58-135
% The figure is descriptive of the method only: no accuracy, no comparison
% between configurations, and no quantity of any kind appears in it.
%
% The picture is a horizontal left-to-right pipeline about 12.9cm wide, so it
% stays inside the report's text width unscaled -- do not add scale= or
% \resizebox. Box labels are intentionally terse: the operational detail lives in
% the prose of \ref{sec:gen-authoring}, \ref{sec:gen-control} and
% \ref{sec:gen-mechanism}, so the figure carries only the dataflow -- the main
% spine (author -> simulate -> validate -> render -> generate), the bounded
% re-author loop dropping below the simulate/validate span, and the separate
% appearance track feeding the hand-off from below.
% =====================================================================
% The TikZ libraries must be loaded in the PREAMBLE, not here: a figure
% environment forms a group, and TikZ defines a library's macros locally, so a
% \usetikzlibrary issued inside a figure does not survive to the next one. This
% guard turns a missing preamble line into a clear message instead of TikZ's
% confusing "You need to say \usetikzlibrary{calc}" further down.
\makeatletter
\@ifundefined{tikz@library@calc@loaded}{%
  \PackageError{fig_authoring_pipeline}{%
    The preamble must contain\MessageBreak
    \string\usetikzlibrary{arrows.meta,positioning,calc}%
  }{Add that line after \string\usepackage{tikz}.}}{}
\makeatother

% The caller's \input runs BEFORE its \caption, so a bare \label here would bind
% to the enclosing section instead of to the figure. Step the figure counter,
% label that value, then hand the counter back: the \caption that follows steps
% it again to the same number, so \cref{fig:authoring} resolves to this figure.
% Outside a float the block does nothing.
\makeatletter
\@ifundefined{@captype}{}{%
  \refstepcounter{figure}%
  \label{fig:authoring}%
  \addtocounter{figure}{-1}%
}
\makeatother

\begingroup
\providecommand{\apsub}[1]{{\scriptsize\color{adobegray}#1}}

\begin{tikzpicture}[
    x=1cm, y=1cm,
    % --- shared box shape: terse title only; detail lives in the prose ----
    apbox/.style   = {rounded corners=2pt, line width=0.7pt, align=center,
                      inner xsep=6pt, inner ysep=4.2pt,
                      font=\sffamily\footnotesize, text=adobeink,
                      minimum height=8.5mm},
    % --- LEARNED component: light gray fill --------------------------------
    aplearn/.style = {apbox, draw=adobeink, fill=adobegray!12},
    % --- DETERMINISTIC component: white with an ink outline ----------------
    apdet/.style   = {apbox, draw=adobeink, fill=white},
    % --- accent: the check a trajectory must pass before it is rendered -----
    apgate/.style  = {apbox, draw=adobered, line width=0.9pt, fill=white,
                      text=adobered},
    % --- what is handed to the system (input) ------------------------------
    apdata/.style  = {apbox, draw=adobegray, fill=white,
                      dash pattern=on 2pt off 1.6pt},
    apline/.style  = {draw=adobeink, line width=0.7pt},
    aparr/.style   = {apline, -{Straight Barb[length=3.6pt,width=3.8pt]}},
    % the re-author loop, drawn in the accent colour so the cycle reads at a glance
    aploop/.style  = {draw=adobered, line width=0.7pt,
                      -{Straight Barb[length=3.6pt,width=3.8pt]}},
    apnote/.style  = {font=\sffamily\scriptsize, text=adobegray,
                      inner sep=0pt, align=left},
    aptag/.style   = {font=\sffamily\scriptsize, text=adobegray, fill=white,
                      inner xsep=3pt, inner ysep=1pt, align=center},
    aptagr/.style  = {aptag, text=adobered},
    apsw/.style    = {rounded corners=1.2pt, line width=0.7pt, inner sep=0pt,
                      minimum width=11pt, minimum height=7.5pt},
  ]

  % Layout: a horizontal main spine along y=0, left to right, sized to the text
  % width (five boxes: author -> simulate -> validate -> render -> generate).
  % The hand-off is folded into the render->generate arrow. The bounded
  % re-author loop drops below the simulate/validate span; the appearance track
  % sits below the render->generate span and joins the hand-off arrow from below.
  \def\bw{1.8cm}          % spine-box width
  \def\hg{0.46}           % horizontal gap between spine boxes
  \tikzset{spine/.style={apbox, text width=\bw, minimum height=13mm,
                         anchor=west}}

  % ---------------- main spine (left -> right) -----------------------------
  \node[spine, apdata] (spec)   at (0,0) {\textbf{Physics request}\\[1pt]\apsub{text}};
  \node[spine, aplearn] (author) at ($(spec.east)+(\hg,0)$)
    {\textbf{LM writes scene spec}};
  \node[spine, apdet] (sim) at ($(author.east)+(\hg,0)$)
    {\textbf{Simulate}\\[1pt]\apsub{renderer's engine}};
  \node[spine, apgate] (gate) at ($(sim.east)+(\hg,0)$)
    {\textbf{Validate trajectory}};
  \node[spine, apdet] (render) at ($(gate.east)+(\hg,0)$)
    {\textbf{Camera + render control}};
  \node[spine, aplearn] (gen) at ($(render.east)+(1.35,0)$)
    {\textbf{Diffusion backbone}};

  % ---------------- re-author loop (accent), below sim/validate ------------
  \node[aplearn, text width=3.4cm, anchor=north, minimum height=8.5mm]
    (repair) at ($(sim.south)!0.5!(gate.south)+(0,-0.92)$)
    {\textbf{Re-author}\ \ \apsub{$\leq 3$ rounds}};

  % ---------------- appearance track (learned), below render->gen ----------
  \node[aplearn, text width=3.7cm, anchor=north, minimum height=8.5mm]
    (texture) at ($(render.south east)!0.5!(gen.south west)+(0,-0.92)$)
    {\textbf{Appearance written separately}\\[1pt]\apsub{no access to the control}};

  % ---------------- edges: main spine --------------------------------------
  \draw[aparr] (spec)    -- (author);
  \draw[aparr] (author)  -- (sim);
  \draw[aparr] (sim)     -- (gate);
  \draw[aparr] (gate)    -- (render)
        node[aptagr, midway, above=0pt] {validated};
  % render -> generate carries the hand-off (control + appearance + seed + strength)
  \draw[aparr] (render)  -- (gen)
        node[aptag, midway, above=0pt] {hand off};

  % ---------------- re-author loop: validate -> re-author -> re-simulate ----
  \draw[aploop] (gate.south) |- (repair.east)
        node[aptagr, pos=0.28, right=1pt] {violated};
  \draw[aploop] (repair.west) -| (sim.south)
        node[aptagr, pos=0.72, left=1pt] {re-simulate};

  % ---------------- appearance joins the hand-off arrow from below ---------
  \draw[aparr] (texture.north) -- ($(render.east)!0.5!(gen.west)$);

  % ---------------- inline legend (below everything) -----------------------
  \coordinate (leg) at ($(spec.west |- repair.south)+(0,-0.62)$);
  \node[apsw, draw=adobeink, fill=adobegray!12, anchor=north west]
    (swA) at (leg) {};
  \node[apnote, right=3pt of swA] (lbA) {learned};

  \node[apsw, draw=adobeink, fill=white, right=9pt of lbA] (swB) {};
  \node[apnote, right=3pt of swB] (lbB) {deterministic};

  \node[apsw, draw=adobered, fill=white, right=9pt of lbB] (swC) {};
  \node[apnote, right=3pt of swC] (lbC) {trajectory check and re-author loop};

  \node[apsw, draw=adobegray, fill=white, dash pattern=on 2pt off 1.6pt,
        right=9pt of lbC] (swD) {};
  \node[apnote, right=3pt of swD] (lbD) {input};

\end{tikzpicture}
\endgroup

%       \caption{...}
%     \end{figure}
% The \label is carried HERE (see the block just below \makeatother), so main.tex
% must NOT add a second \label{fig:authoring}.
%
% CONTENT follows the subsection "From a physics description to a control video"
% (\ref{sec:gen-authoring}) and is drawn from the implementation in
% ~/diffphy_exp014/scripts/exp014/agentic_mask.py -- stage order from build_one()
% at lines 589-663, with:
%   text-in, no reference video   docstring l.4-8; author() l.294-306; main() l.683
%   scene specification written   author() l.294-327; scene contract l.124-149
%   required fields re-asked once _validate() l.249-265; author() l.322-326
%   waypoints only where needed   docstring l.26-29; motion model l.98-117
%   simulated by the renderer's   _sim_traj() l.348-352 (expand + simulate);
%     own engine                  silhouette_render.render_frames l.209 calls the
%                                 same simulate(); physics_world.py is pymunk
%   validation on the trajectory  lint_scene() l.355-450 (ballistic must be
%                                 dynamic l.406; off-frame l.410-417; rest /
%                                 support l.419-424; rising must rise l.426-428;
%                                 reach / miss targets l.430-445)
%   repair loop, three rounds     build_one() l.604-612; repair() l.472-488;
%                                 realized-motion readout _traj_summary l.453-469
%   appearance written apart      author_texture() l.272-291; called l.616
%   camera + graphic-body strip   apply_camera() l.491-504; _strip_graphic_bodies
%                                 l.210-215; called l.618-619
%   render + contact sheet        build_one l.623-627; render_gif; strip l.507-521
%   what the generator receives   build_one l.632-662 (fill-ratio cap l.632-638,
%                                 low-control-strength cap l.652-654), consumed by
%                                 generate_worker.py l.176-183 -> video_stylizer
%                                 WanVACEStylizer.stylize l.58-135
% The figure is descriptive of the method only: no accuracy, no comparison
% between configurations, and no quantity of any kind appears in it.
%
% The picture is a horizontal left-to-right pipeline about 12.9cm wide, so it
% stays inside the report's text width unscaled -- do not add scale= or
% \resizebox. Box labels are intentionally terse: the operational detail lives in
% the prose of \ref{sec:gen-authoring}, \ref{sec:gen-control} and
% \ref{sec:gen-mechanism}, so the figure carries only the dataflow -- the main
% spine (author -> simulate -> validate -> render -> generate), the bounded
% re-author loop dropping below the simulate/validate span, and the separate
% appearance track feeding the hand-off from below.
% =====================================================================
% The TikZ libraries must be loaded in the PREAMBLE, not here: a figure
% environment forms a group, and TikZ defines a library's macros locally, so a
% \usetikzlibrary issued inside a figure does not survive to the next one. This
% guard turns a missing preamble line into a clear message instead of TikZ's
% confusing "You need to say \usetikzlibrary{calc}" further down.
\makeatletter
\@ifundefined{tikz@library@calc@loaded}{%
  \PackageError{fig_authoring_pipeline}{%
    The preamble must contain\MessageBreak
    \string\usetikzlibrary{arrows.meta,positioning,calc}%
  }{Add that line after \string\usepackage{tikz}.}}{}
\makeatother

% The caller's \input runs BEFORE its \caption, so a bare \label here would bind
% to the enclosing section instead of to the figure. Step the figure counter,
% label that value, then hand the counter back: the \caption that follows steps
% it again to the same number, so \cref{fig:authoring} resolves to this figure.
% Outside a float the block does nothing.
\makeatletter
\@ifundefined{@captype}{}{%
  \refstepcounter{figure}%
  \label{fig:authoring}%
  \addtocounter{figure}{-1}%
}
\makeatother

\begingroup
\providecommand{\apsub}[1]{{\scriptsize\color{adobegray}#1}}

\begin{tikzpicture}[
    x=1cm, y=1cm,
    % --- shared box shape: terse title only; detail lives in the prose ----
    apbox/.style   = {rounded corners=2pt, line width=0.7pt, align=center,
                      inner xsep=6pt, inner ysep=4.2pt,
                      font=\sffamily\footnotesize, text=adobeink,
                      minimum height=8.5mm},
    % --- LEARNED component: light gray fill --------------------------------
    aplearn/.style = {apbox, draw=adobeink, fill=adobegray!12},
    % --- DETERMINISTIC component: white with an ink outline ----------------
    apdet/.style   = {apbox, draw=adobeink, fill=white},
    % --- accent: the check a trajectory must pass before it is rendered -----
    apgate/.style  = {apbox, draw=adobered, line width=0.9pt, fill=white,
                      text=adobered},
    % --- what is handed to the system (input) ------------------------------
    apdata/.style  = {apbox, draw=adobegray, fill=white,
                      dash pattern=on 2pt off 1.6pt},
    apline/.style  = {draw=adobeink, line width=0.7pt},
    aparr/.style   = {apline, -{Straight Barb[length=3.6pt,width=3.8pt]}},
    % the re-author loop, drawn in the accent colour so the cycle reads at a glance
    aploop/.style  = {draw=adobered, line width=0.7pt,
                      -{Straight Barb[length=3.6pt,width=3.8pt]}},
    apnote/.style  = {font=\sffamily\scriptsize, text=adobegray,
                      inner sep=0pt, align=left},
    aptag/.style   = {font=\sffamily\scriptsize, text=adobegray, fill=white,
                      inner xsep=3pt, inner ysep=1pt, align=center},
    aptagr/.style  = {aptag, text=adobered},
    apsw/.style    = {rounded corners=1.2pt, line width=0.7pt, inner sep=0pt,
                      minimum width=11pt, minimum height=7.5pt},
  ]

  % Layout: a horizontal main spine along y=0, left to right, sized to the text
  % width (five boxes: author -> simulate -> validate -> render -> generate).
  % The hand-off is folded into the render->generate arrow. The bounded
  % re-author loop drops below the simulate/validate span; the appearance track
  % sits below the render->generate span and joins the hand-off arrow from below.
  \def\bw{1.8cm}          % spine-box width
  \def\hg{0.46}           % horizontal gap between spine boxes
  \tikzset{spine/.style={apbox, text width=\bw, minimum height=13mm,
                         anchor=west}}

  % ---------------- main spine (left -> right) -----------------------------
  \node[spine, apdata] (spec)   at (0,0) {\textbf{Physics request}\\[1pt]\apsub{text}};
  \node[spine, aplearn] (author) at ($(spec.east)+(\hg,0)$)
    {\textbf{LM writes scene spec}};
  \node[spine, apdet] (sim) at ($(author.east)+(\hg,0)$)
    {\textbf{Simulate}\\[1pt]\apsub{renderer's engine}};
  \node[spine, apgate] (gate) at ($(sim.east)+(\hg,0)$)
    {\textbf{Validate trajectory}};
  \node[spine, apdet] (render) at ($(gate.east)+(\hg,0)$)
    {\textbf{Camera + render control}};
  \node[spine, aplearn] (gen) at ($(render.east)+(1.35,0)$)
    {\textbf{Diffusion backbone}};

  % ---------------- re-author loop (accent), below sim/validate ------------
  \node[aplearn, text width=3.4cm, anchor=north, minimum height=8.5mm]
    (repair) at ($(sim.south)!0.5!(gate.south)+(0,-0.92)$)
    {\textbf{Re-author}\ \ \apsub{$\leq 3$ rounds}};

  % ---------------- appearance track (learned), below render->gen ----------
  \node[aplearn, text width=3.7cm, anchor=north, minimum height=8.5mm]
    (texture) at ($(render.south east)!0.5!(gen.south west)+(0,-0.92)$)
    {\textbf{Appearance written separately}\\[1pt]\apsub{no access to the control}};

  % ---------------- edges: main spine --------------------------------------
  \draw[aparr] (spec)    -- (author);
  \draw[aparr] (author)  -- (sim);
  \draw[aparr] (sim)     -- (gate);
  \draw[aparr] (gate)    -- (render)
        node[aptagr, midway, above=0pt] {validated};
  % render -> generate carries the hand-off (control + appearance + seed + strength)
  \draw[aparr] (render)  -- (gen)
        node[aptag, midway, above=0pt] {hand off};

  % ---------------- re-author loop: validate -> re-author -> re-simulate ----
  \draw[aploop] (gate.south) |- (repair.east)
        node[aptagr, pos=0.28, right=1pt] {violated};
  \draw[aploop] (repair.west) -| (sim.south)
        node[aptagr, pos=0.72, left=1pt] {re-simulate};

  % ---------------- appearance joins the hand-off arrow from below ---------
  \draw[aparr] (texture.north) -- ($(render.east)!0.5!(gen.west)$);

  % ---------------- inline legend (below everything) -----------------------
  \coordinate (leg) at ($(spec.west |- repair.south)+(0,-0.62)$);
  \node[apsw, draw=adobeink, fill=adobegray!12, anchor=north west]
    (swA) at (leg) {};
  \node[apnote, right=3pt of swA] (lbA) {learned};

  \node[apsw, draw=adobeink, fill=white, right=9pt of lbA] (swB) {};
  \node[apnote, right=3pt of swB] (lbB) {deterministic};

  \node[apsw, draw=adobered, fill=white, right=9pt of lbB] (swC) {};
  \node[apnote, right=3pt of swC] (lbC) {trajectory check and re-author loop};

  \node[apsw, draw=adobegray, fill=white, dash pattern=on 2pt off 1.6pt,
        right=9pt of lbC] (swD) {};
  \node[apnote, right=3pt of swD] (lbD) {input};

\end{tikzpicture}
\endgroup

%       \caption{...}
%     \end{figure}
% The \label is carried HERE (see the block just below \makeatother), so main.tex
% must NOT add a second \label{fig:authoring}.
%
% CONTENT follows the subsection "From a physics description to a control video"
% (\ref{sec:gen-authoring}) and is drawn from the implementation in
% ~/diffphy_exp014/scripts/exp014/agentic_mask.py -- stage order from build_one()
% at lines 589-663, with:
%   text-in, no reference video   docstring l.4-8; author() l.294-306; main() l.683
%   scene specification written   author() l.294-327; scene contract l.124-149
%   required fields re-asked once _validate() l.249-265; author() l.322-326
%   waypoints only where needed   docstring l.26-29; motion model l.98-117
%   simulated by the renderer's   _sim_traj() l.348-352 (expand + simulate);
%     own engine                  silhouette_render.render_frames l.209 calls the
%                                 same simulate(); physics_world.py is pymunk
%   validation on the trajectory  lint_scene() l.355-450 (ballistic must be
%                                 dynamic l.406; off-frame l.410-417; rest /
%                                 support l.419-424; rising must rise l.426-428;
%                                 reach / miss targets l.430-445)
%   repair loop, three rounds     build_one() l.604-612; repair() l.472-488;
%                                 realized-motion readout _traj_summary l.453-469
%   appearance written apart      author_texture() l.272-291; called l.616
%   camera + graphic-body strip   apply_camera() l.491-504; _strip_graphic_bodies
%                                 l.210-215; called l.618-619
%   render + contact sheet        build_one l.623-627; render_gif; strip l.507-521
%   what the generator receives   build_one l.632-662 (fill-ratio cap l.632-638,
%                                 low-control-strength cap l.652-654), consumed by
%                                 generate_worker.py l.176-183 -> video_stylizer
%                                 WanVACEStylizer.stylize l.58-135
% The figure is descriptive of the method only: no accuracy, no comparison
% between configurations, and no quantity of any kind appears in it.
%
% The picture is a horizontal left-to-right pipeline about 12.9cm wide, so it
% stays inside the report's text width unscaled -- do not add scale= or
% \resizebox. Box labels are intentionally terse: the operational detail lives in
% the prose of \ref{sec:gen-authoring}, \ref{sec:gen-control} and
% \ref{sec:gen-mechanism}, so the figure carries only the dataflow -- the main
% spine (author -> simulate -> validate -> render -> generate), the bounded
% re-author loop dropping below the simulate/validate span, and the separate
% appearance track feeding the hand-off from below.
% =====================================================================
% The TikZ libraries must be loaded in the PREAMBLE, not here: a figure
% environment forms a group, and TikZ defines a library's macros locally, so a
% \usetikzlibrary issued inside a figure does not survive to the next one. This
% guard turns a missing preamble line into a clear message instead of TikZ's
% confusing "You need to say \usetikzlibrary{calc}" further down.
\makeatletter
\@ifundefined{tikz@library@calc@loaded}{%
  \PackageError{fig_authoring_pipeline}{%
    The preamble must contain\MessageBreak
    \string\usetikzlibrary{arrows.meta,positioning,calc}%
  }{Add that line after \string\usepackage{tikz}.}}{}
\makeatother

% The caller's \input runs BEFORE its \caption, so a bare \label here would bind
% to the enclosing section instead of to the figure. Step the figure counter,
% label that value, then hand the counter back: the \caption that follows steps
% it again to the same number, so \cref{fig:authoring} resolves to this figure.
% Outside a float the block does nothing.
\makeatletter
\@ifundefined{@captype}{}{%
  \refstepcounter{figure}%
  \label{fig:authoring}%
  \addtocounter{figure}{-1}%
}
\makeatother

\begingroup
\providecommand{\apsub}[1]{{\scriptsize\color{adobegray}#1}}

\begin{tikzpicture}[
    x=1cm, y=1cm,
    % --- shared box shape: terse title only; detail lives in the prose ----
    apbox/.style   = {rounded corners=2pt, line width=0.7pt, align=center,
                      inner xsep=6pt, inner ysep=4.2pt,
                      font=\sffamily\footnotesize, text=adobeink,
                      minimum height=8.5mm},
    % --- LEARNED component: light gray fill --------------------------------
    aplearn/.style = {apbox, draw=adobeink, fill=adobegray!12},
    % --- DETERMINISTIC component: white with an ink outline ----------------
    apdet/.style   = {apbox, draw=adobeink, fill=white},
    % --- accent: the check a trajectory must pass before it is rendered -----
    apgate/.style  = {apbox, draw=adobered, line width=0.9pt, fill=white,
                      text=adobered},
    % --- what is handed to the system (input) ------------------------------
    apdata/.style  = {apbox, draw=adobegray, fill=white,
                      dash pattern=on 2pt off 1.6pt},
    apline/.style  = {draw=adobeink, line width=0.7pt},
    aparr/.style   = {apline, -{Straight Barb[length=3.6pt,width=3.8pt]}},
    % the re-author loop, drawn in the accent colour so the cycle reads at a glance
    aploop/.style  = {draw=adobered, line width=0.7pt,
                      -{Straight Barb[length=3.6pt,width=3.8pt]}},
    apnote/.style  = {font=\sffamily\scriptsize, text=adobegray,
                      inner sep=0pt, align=left},
    aptag/.style   = {font=\sffamily\scriptsize, text=adobegray, fill=white,
                      inner xsep=3pt, inner ysep=1pt, align=center},
    aptagr/.style  = {aptag, text=adobered},
    apsw/.style    = {rounded corners=1.2pt, line width=0.7pt, inner sep=0pt,
                      minimum width=11pt, minimum height=7.5pt},
  ]

  % Layout: a horizontal main spine along y=0, left to right, sized to the text
  % width (five boxes: author -> simulate -> validate -> render -> generate).
  % The hand-off is folded into the render->generate arrow. The bounded
  % re-author loop drops below the simulate/validate span; the appearance track
  % sits below the render->generate span and joins the hand-off arrow from below.
  \def\bw{1.8cm}          % spine-box width
  \def\hg{0.46}           % horizontal gap between spine boxes
  \tikzset{spine/.style={apbox, text width=\bw, minimum height=13mm,
                         anchor=west}}

  % ---------------- main spine (left -> right) -----------------------------
  \node[spine, apdata] (spec)   at (0,0) {\textbf{Physics request}\\[1pt]\apsub{text}};
  \node[spine, aplearn] (author) at ($(spec.east)+(\hg,0)$)
    {\textbf{LM writes scene spec}};
  \node[spine, apdet] (sim) at ($(author.east)+(\hg,0)$)
    {\textbf{Simulate}\\[1pt]\apsub{renderer's engine}};
  \node[spine, apgate] (gate) at ($(sim.east)+(\hg,0)$)
    {\textbf{Validate trajectory}};
  \node[spine, apdet] (render) at ($(gate.east)+(\hg,0)$)
    {\textbf{Camera + render control}};
  \node[spine, aplearn] (gen) at ($(render.east)+(1.35,0)$)
    {\textbf{Diffusion backbone}};

  % ---------------- re-author loop (accent), below sim/validate ------------
  \node[aplearn, text width=3.4cm, anchor=north, minimum height=8.5mm]
    (repair) at ($(sim.south)!0.5!(gate.south)+(0,-0.92)$)
    {\textbf{Re-author}\ \ \apsub{$\leq 3$ rounds}};

  % ---------------- appearance track (learned), below render->gen ----------
  \node[aplearn, text width=3.7cm, anchor=north, minimum height=8.5mm]
    (texture) at ($(render.south east)!0.5!(gen.south west)+(0,-0.92)$)
    {\textbf{Appearance written separately}\\[1pt]\apsub{no access to the control}};

  % ---------------- edges: main spine --------------------------------------
  \draw[aparr] (spec)    -- (author);
  \draw[aparr] (author)  -- (sim);
  \draw[aparr] (sim)     -- (gate);
  \draw[aparr] (gate)    -- (render)
        node[aptagr, midway, above=0pt] {validated};
  % render -> generate carries the hand-off (control + appearance + seed + strength)
  \draw[aparr] (render)  -- (gen)
        node[aptag, midway, above=0pt] {hand off};

  % ---------------- re-author loop: validate -> re-author -> re-simulate ----
  \draw[aploop] (gate.south) |- (repair.east)
        node[aptagr, pos=0.28, right=1pt] {violated};
  \draw[aploop] (repair.west) -| (sim.south)
        node[aptagr, pos=0.72, left=1pt] {re-simulate};

  % ---------------- appearance joins the hand-off arrow from below ---------
  \draw[aparr] (texture.north) -- ($(render.east)!0.5!(gen.west)$);

  % ---------------- inline legend (below everything) -----------------------
  \coordinate (leg) at ($(spec.west |- repair.south)+(0,-0.62)$);
  \node[apsw, draw=adobeink, fill=adobegray!12, anchor=north west]
    (swA) at (leg) {};
  \node[apnote, right=3pt of swA] (lbA) {learned};

  \node[apsw, draw=adobeink, fill=white, right=9pt of lbA] (swB) {};
  \node[apnote, right=3pt of swB] (lbB) {deterministic};

  \node[apsw, draw=adobered, fill=white, right=9pt of lbB] (swC) {};
  \node[apnote, right=3pt of swC] (lbC) {trajectory check and re-author loop};

  \node[apsw, draw=adobegray, fill=white, dash pattern=on 2pt off 1.6pt,
        right=9pt of lbC] (swD) {};
  \node[apnote, right=3pt of swD] (lbD) {input};

\end{tikzpicture}
\endgroup

  \caption{Physics-guided generation. A \emph{physics request} becomes a scene
  that is simulated and validated, revising on failure (\textcolor{adobered}{red}~loop);
  the validated motion is a depth control that, with an \emph{appearance prompt},
  conditions the generator. Right: depth control and generated clip at matched
  timestamps, tracing the ball's path.}
  \label{fig:genpipe}
\end{figure}

\subsection{Preliminaries}
\label{sec:gen-prelim}

The atomic unit is a \emph{clip--prompt pair}: a clip
$\mathcal{V}=\{(I_t,\tau_t)\}_{t\in\mathcal{T}_{\mathcal V}}$, a time-ordered sequence of frames $I_t$
at timestamps $\tau_t$ with duration $D_{\mathcal V}$ (\cref{sec:critic}), and a natural-language
prompt $p$. The prompt supplies appearance only (subject, material, setting), the geometric control
$C$ carrying the physical target---a solver-rendered video with per-frame per-pixel depth intensity
$C_t(u)$ (\cref{sec:gen-control}); $\hat{\mathcal{V}}$ denotes a generated sample, $\mathcal{V}$ any
clip fed to the critic. Generation and evaluation are thus two maps over a clip and its prompt, built
and benchmarked separately,
\begin{equation}
  \hat{\mathcal{V}}\sim\operatorname{Gen}(p,C)
  \qquad\qquad
  Y=\operatorname{Crit}(\mathcal{V},p) ,
  \label{eq:gen-crit-maps}
\end{equation}
where $\operatorname{Gen}$ samples a clip conditioned on $p$, $C$, a mask $M$ and a scalar strength
$s$---asserting no physical fidelity---and $\operatorname{Crit}$ is the auditable critic \PhyReAct{}
of \cref{sec:critic}, returning the clip-level verdict $Y$ (\cref{eq:claim-clip-rollup}).

\subsection{From specification to registered control}\label{sec:gen-control}

This stage produces the geometric control, validating the motion it carries before any control is
rendered, so the generator is conditioned on a verified trajectory rather than an unconstrained
request. It is a three-map chain from the event description $p_{\mathrm{ev}}$ to the control $C$,
\begin{equation}
p_{\mathrm{ev}}
\;\xrightarrow{\;\textsc{Author}\;}\; \mathcal{S}
\;\xrightarrow{\;\textsc{Sim}\;}\; \tau
\;\xrightarrow{\;\textsc{Render}(\,\cdot\,;\,\mathcal{S})\;}\; C,
\label{eq:control-chain}
\end{equation}
where $\textsc{Author}$ has a language model compile the text into a structured scene specification
$\mathcal{S}$ (bodies, initial states, colliders, static surfaces, and a camera; \cref{fig:authoring}),
$\textsc{Sim}$ runs the physics solver to a trajectory $\tau$, and $\textsc{Render}$ converts the
solver state into the depth control $C$ of \cref{sec:gen-control}. The middle map is gated by a
validation loop (\cref{eq:scene-validate}) so only a $\tau$ meeting the event's physical predicates is
rendered; the three maps are developed in turn below.

A body is simulated from its initial state when the engine determines the motion (a ballistic arc and
its contacts), while prescribed waypoints are reserved for externally actuated or explicitly
counterfactual motion. Validation applies to the simulated trajectory, not the specification text,
through a conjunction of event-specific predicates $\Phi_{\mathrm{ev}}$: a ballistic segment must form
an arc, satisfy the required hit or miss relation, settle when requested, and remain visible. The
$\textsc{Sim}$ map is therefore a bounded revise-and-resimulate loop, $\operatorname{Solve}$ running
the solver to a trajectory and $\textsc{Summarize}$ reducing it to a text report of its realized
motion. Starting from $\mathcal{S}^{(0)}=\mathcal{S}$, for round $k=0,1,2$,
\begin{equation}
\begin{aligned}
\tau^{(k)} &= \operatorname{Solve}\!\big(\mathcal{S}^{(k)}\big),\\
\mathcal{S}^{(k+1)} &= \textsc{Author}\!\big(\mathcal{S}^{(k)},\, \textsc{Summarize}(\tau^{(k)})\big),
\end{aligned}
\label{eq:scene-validate}
\end{equation}
where the revision on the second line is taken only while $\neg\,\Phi_{\mathrm{ev}}(\tau^{(k)})$: a
failed predicate returns a summary of the realized motion to the authoring model, which revises the
specification, for at most three revisions (four simulations). The first $\tau^{(k)}$ with
$\Phi_{\mathrm{ev}}(\tau^{(k)})$ true is the validated trajectory $\tau$ of \cref{eq:control-chain},
passed to $\textsc{Render}$; exhausting the budget yields an explicitly marked unvalidated fallback
rather than a validated scene.

The $\textsc{Render}$ map turns the validated trajectory $\tau$ and its scene $\mathcal{S}$ into the
control $C$. Rendered from solver state rather than inferred from a generated video, $C$ is registered
to the intended motion by construction, all streams sharing the scene's frame clock (solver, camera,
and playback settings deferred to \cref{sec:gen-setup}). For each of the clip's $T$ frames, rendering
places $\mathcal{S}$'s geometry at the per-frame pose, images it with the scene camera, and reads a
depth buffer $D_t(u)$---camera distance at pixel $u$ in frame $t$---normalized to a near-bright,
far-dark intensity $C_t(u)=1-\operatorname{clip}\big((D_t(u)-d_-)/(d_+-d_-),0,1\big)\in[0,1]$, where
$d_-,d_+$ are the $1$st and $97$th percentiles of the non-background depths pooled over the whole clip
and $\operatorname{clip}$ clamps to $[0,1]$; the control $C=(C_t)_{t=0}^{T-1}$ then conditions the
generator in \cref{sec:gen-mechanism}. Background ($D_t(u)$
beyond $d_+$) maps to $0$; pooling the percentiles once fixes the mapping so the same distance is the
same intensity throughout, the cut asymmetric because the far tail abuts the background and is
dominated by depth-edge noise, with a small Gaussian blur softening rasterized boundaries. When
$d_+>d_-$ fails---no non-background pixels, coincident percentiles, or a predeclared flat
scene---$\textsc{Render}$ falls back to a binary moving-body silhouette, depth otherwise preferred for
preserving both outline and distance ordering.

\subsection{Physics-conditioned video generation}\label{sec:gen-mechanism}

\paragraph{The backbone.} The generator is a latent video-diffusion transformer that samples a clip
by iteratively denoising a latent from noise. We use Wan~2.2~\cite{wan2025}: a DiT-style denoiser run
along a flow-matching schedule of $S$ UniPC multistep iterations, indexed by $i$ from high noise to
low as $\sigma_i$ decreases (\cref{eq:flow-sigma}). On its own it is a text-to-video generator
conditioned only on an appearance prompt $p$ (subject, material, setting), with no input through which
a specified motion can be imposed.

\paragraph{Adding a control branch (VACE).} To make the motion controllable we adopt
VACE~\cite{vace2025}, which augments a frozen backbone with a \emph{control branch}: a control video
and a generation mask are encoded into control features added into the backbone's transformer blocks,
so an external spatiotemporal signal steers generation without retraining the base
weights~\cite{zhang2023controlnet,wang2023videocomposer}. The composite Wan~2.2-VACE takes four
inputs: the appearance prompt $p$, a geometric control video $C$, a generation mask $M$, and a scalar
conditioning strength $s$.

\paragraph{Fusing the simulated control.} We drive VACE's control branch with the registered depth
control of \cref{sec:gen-control}: $C$ is the depth video $(C_t)_{t=0}^{T-1}$ carrying the validated
motion, and the mask decides which pixels VACE may synthesize. We set $M$ all-white, so every output
pixel is generated and no pixel of the simulator render is copied in---the control acts only through
the branch, never by pasting appearance. The branch then injects, at each transformer layer $\ell$ and
step $i$, an additive update to the hidden state $h_i^{(\ell)}$; with $z^{\mathrm{ctrl}}$ the encoded
control and $W^{(\ell)}$ the branch's projection at layer $\ell$,
\begin{equation}
h_i^{(\ell)} \;\leftarrow\; h_i^{(\ell)} \;+\; s\,W^{(\ell)}z^{\mathrm{ctrl}},
\label{eq:control-release}
\end{equation}
so the control enters as a strength-$s$ residual rather than by overwriting pixels. The same $s$ is
applied at every control layer and stays active throughout sampling in the deployed generator, lowered
when the control covers little of the frame. Conditioned on $(p,C,M,s)$ with a fixed seed and a set of
excluded terms, the backbone induces a distribution over clips from which
$\hat{\mathcal{V}}\sim\operatorname{Gen}(p,C)$ is drawn---a sample, asserting no physical fidelity.

\subsection{Control withdrawal as a mechanism probe}\label{sec:gen-withdrawal}

\emph{When} can the control be released? The injection of \cref{eq:control-release} runs throughout
sampling, but if the trajectory settles by some intermediate noise level the control is redundant
afterward. Switching it off partway is awkward for VACE's additive residual, so we study the question
on a \emph{separate} Wan2.2-derived backbone whose control path is easy to interrupt---a mechanism
probe, not the deployed generator---and read the commitment point back as a property of the Wan~2.2
diffusion prior both share. Run at $40$ UniPC steps with flow shift $5.0$, this backbone has no
trained control branch: its only conditioning path encodes a rendered source video into context
tokens concatenated with the text and image tokens, so withdrawal is directly available---the source
tokens are removed from the attended sequence at a chosen step, distinct from the additive-residual
injection of \cref{eq:control-release}: over $S$ steps the source tokens are present at step $i$ iff
$i<i^{*}$, where $i^{*}=\operatorname{clip}(\operatorname{round}(fS),0,S-1)$ and $f\in[0,1]$ is the
withdrawal fraction. Steps before $i^{*}$ attend over the source tokens at the scale
$\omega$ the deployed recipe uses; from $i^{*}$ onward the sequence is identical to that of an
unconditioned generation.

Because $i^{*}$ is an index, its noise level depends on the schedule, so the withdrawal point is
reported as $\sigma_{\mathrm{rel}}=\sigma_{i^{*}}$, the level the flow-matching schedule assigns step
$i^{*}$ (the shift-dependent map is \cref{eq:flow-sigma} in \cref{sec:app-schedule}); two runs at the
same step index thus hit different noise levels whenever $S$ or shift $\kappa$ differs, which is why
\cref{sec:gen-release} indexes by $\sigma_{\mathrm{rel}}$. Finally, the multistep solver forms each
update from derivatives cached at previous steps, which after a withdrawal still carry the removed
tokens' influence; the measurement compares one variant that clears that history at $i^{*}$ against
one that keeps it, otherwise identical.
